\section{Decision-level sensitivity of estimated constraints}
\label{app:plugin-v12}
Appendix~\ref{app:proofs} bounds score perturbations. Moving from scores to
decisions additionally requires stability of the feasible set. Constraint
tightening under strict feasibility is established in constrained learning
\citep[Proposition 4.1]{rigollet2011strict}; consistent plug-in methods also
exist for general confusion-matrix objectives and constraints
\citep{tavker2020plugin}. The result below specializes this sensitivity
argument to accepted exposure and ratio risk. We do not claim a new general
plug-in framework or use the result to certify our empirical policies.

Fix a distribution and nonnegative integrable conditional moments $\ell,w$
with $T=\E w>0$. For a gate $a$, write
$M(a)=\E[aw]$ and $G(a)=\E[a(\ell-rw)]$, where $r\geq0$.
Use hats for the same expectations with estimated moments. A common convex
class $\mathcal A\subseteq\{a:0\leq a\leq1\}$ may include the true case floor.
Suppose
\[
 \sup_{a\in\mathcal A}|\widehat M(a)-M(a)|\leq\epsilon_M,
 \qquad
 \sup_{a\in\mathcal A}|\widehat G(a)-G(a)|\leq\epsilon_G.
\]
Sufficient choices are $\epsilon_M=\E|\widehat w-w|$ and
$\epsilon_G=\E|\widehat\ell-\ell|+r\E|\widehat w-w|$.

\begin{proposition}[Buffered plug-in sensitivity]
\label{prop:plugin-v12}
Assume a maximizer $a^*$ of $M(a)$ subject to $a\in\mathcal A$, $G(a)\leq0$
exists, and some $a_0\in\mathcal A$ satisfies $G(a_0)\leq-s<0$.
Let $\widehat a$ be an $\eta$-optimal solution maximizing $\widehat M$
subject to $\widehat G\leq-\tau$, where $\eta,\tau\geq0$ and
$\epsilon_G+\tau\leq s$. Then
\begin{align}
 d_W(a^*)-d_W(\widehat a)
 &\leq\frac{\epsilon_G+\tau}{s}
       +\frac{2\epsilon_M+\eta}{T},\label{eq:plugin-cover-v12}\\
 G(\widehat a)&\leq\epsilon_G-\tau.\label{eq:plugin-excess-v12}
\end{align}
If $M(\widehat a)\geq m_0>0$, then
$[R(\widehat a)-r]_+\leq(\epsilon_G-\tau)_+/m_0$.
In particular, $\tau=\epsilon_G$ ensures true feasibility when
$2\epsilon_G\leq s$ and accepted exposure is positive.
\end{proposition}

\begin{proof}
Set $t=(\epsilon_G+\tau)/s\in[0,1]$ and
$a_t=(1-t)a^*+ta_0\in\mathcal A$. Linearity gives
$G(a_t)\leq-ts$, hence $\widehat G(a_t)\leq-\tau$.
Optimality up to $\eta$ and the objective error bound give
\[
 M(\widehat a)\geq\widehat M(\widehat a)-\epsilon_M
 \geq\widehat M(a_t)-\eta-\epsilon_M
 \geq M(a_t)-2\epsilon_M-\eta.
\]
Thus $M(a^*)-M(\widehat a)\leq
t\{M(a^*)-M(a_0)\}+2\epsilon_M+\eta$.
Since $a_0$ is feasible, $0\leq M(a^*)-M(a_0)\leq T$.
Equation~\eqref{eq:plugin-excess-v12} follows from feasibility of
$\widehat a$ for the estimated problem and the constraint error bound.
Dividing by positive $M(\widehat a)$ proves the risk statement.
\end{proof}

The unbuffered output can be truly infeasible, so the coverage shortfall in
Equation~\eqref{eq:plugin-cover-v12} need not be nonnegative. A case floor
alone does not give $m_0$ if $w$ approaches zero. The stronger condition
$w\geq u>0$, $c(a)\geq c_0>0$ gives $m_0=uc_0$.

\paragraph{A finite-menu corollary and its scope.}
For a common finite class, convexity is unnecessary if a true comparator
$a^*$ satisfies $G(a^*)\leq-(\epsilon_G+\tau)$: it remains feasible for
the estimated problem, and the same comparison yields
$M(a^*)-M(\widehat a)\leq2\epsilon_M+\eta$.
The empirical procedure uses fitted heads to propose rankings, but observed
calibration losses to set thresholds. Its five candidates are not closed
under mixtures. Applying either result to that procedure would require
uniform errors for the relevant data-dependent criteria, robustness of
candidate eligibility, and (for oracle comparison) class approximation
control. These assumptions are not supplied by the reported bootstrap
intervals. Changing fitted heads also changes the candidate class.

\paragraph{Why head accuracy alone is insufficient.}
Let two contexts have equal probability, $w=1$ and $\ell=r>0$ in both.
Accept-all is truly feasible. Set $\widehat w=1$ and
$\widehat\ell=(r+\varepsilon,r-\varepsilon^2)$ with
$0<\varepsilon<\min(1,\sqrt r)$. The plug-in optimum accepts the second
context and an $\varepsilon$ fraction of the first, since its constraint
is $\varepsilon a_1\leq\varepsilon^2a_2$.
Its coverage is $(1+\varepsilon)/2$, so its shortfall tends to $1/2$
despite vanishing uniform head error and accepted exposure at least $1/2$.
It meets a .35 case floor and has true risk exactly $r$. No gate has strict
true slack. Hence head error and a positive denominator alone cannot imply
vanishing coverage regret.

For a separate denominator example, let two equally likely contexts have
$w=(2\varepsilon,2(1-\varepsilon))$ and
$\ell=(2(r+1)\varepsilon,2(r+2)(1-\varepsilon))$.
Estimate only the first loss as $2r\varepsilon$. The plug-in optimum
accepts only that context. The loss-head $L^1$ error is $\varepsilon$,
total exposure is one, and case coverage is $1/2$, yet accepted exposure
is $\varepsilon$ and true risk is $r+1$. This example concerns risk
certification; it does not posit a true positive-exposure feasible gate.

Neither example establishes why Descending beats Ratio in M5. Moment errors
can cancel in the ratio, and empirical eligibility and time shift introduce
additional effects. The new result identifies sufficient conditions and
limits of a plug-in guarantee; the empirical estimation mechanism remains
unidentified.
